% part: intuitionistic-logic
% chapter: propositions-as-types
% section: terms-to-proofs

\documentclass[../../../include/open-logic-section]{subfiles}

\begin{document}

\olfileid{int}{pty}{tp}

\olsection{بازیابی \usetoken{P}{derivation} از ترم‌های برهان}

اکنون جهت دیگر را در نظر بگیریم: ترجمهٔ یک ترم برهان درست و نوع‌پذیر
~$N$ به !!a{proof} در~\Log{N2i}. این بازسازی تنها نسبت به بافتی ممکن
است که برای هر متغیر~$x$ آزاد در ترم برهان داده‌شده، جفتی به‌شکل
$x:!A$ در بر داشته باشد. یک ترم خام که در چنین بافتی نوع‌پذیر نباشد
لزوماً هیچ برهانی را رمزگذاری نمی‌کند. با نمونه‌ای بدون متغیر آزاد آغاز
می‌کنیم، یعنی ترم
\[
\lambd[\typeof{x}{!A \land
  !B}][\pair{\proj{2}{x}}{\proj{1}{x}}]
\]
این ترم به شکل~$\lambd[\typeof{x}{!A \land !B}][N_1]$ است. چون تنها
قاعدهٔ~\Log{tN2i} که ترمی از این شکل می‌سازد~$\Intro{\lif}$ است،
می‌دانیم !!{proof} رمزگذاری‌شده به‌وسیلهٔ~$N$ باید با~$\Intro{\lif}$
پایان یابد؛ یعنی پایان آن چنین است:
  \[
  \Axiom$x : !A\land !B \fCenter \pair{\proj{2}{x}}{\proj{1}{x}} :
    !C$
  \RightLabel{\Intro\lif}
  \UnaryInf$\fCenter \lambd[\typeof{x}{!A \land
  !B}][\pair{\proj{2}{x}}{\proj{1}{x}}] : (!A \land !B) \lif !C$
  \DisplayProof
  \]
هنوز نمی‌دانیم~$!C$ چیست، اما می‌دانیم مقدم شرطی~$!A \land !B$ است
و بافت مقدمه باید~$x:!A \land !B$ را در بر داشته باشد.

اکنون ترم برهان متناظر با مقدمه از پیش تعیین شده است:
$\pair{\proj{2}{x}}{\proj{1}{x}}$. !!{proof}ی که این ترم رمزگذاری
می‌کند باید با~$\Intro{\land}$ پایان یابد. هنوز نمی‌دانیم~$!C$ چیست،
اما چون نتیجهٔ~$\Intro{\land}$ است، باید عطف باشد؛ یعنی
$!C \ident (!C_1 \land !C_2)$. بازسازی را ادامه می‌دهیم:
  \[
  \Axiom$x : !A\land !B \fCenter \proj{2}{x} : !C_1$
  \Axiom$x : !A\land !B \fCenter \proj{1}{x} : !C_2$
  \RightLabel{\Intro\land}
  \BinaryInf$x : !A\land !B \fCenter \pair{\proj{2}{x}}{\proj{1}{x}} :
    !C_1 \land !C_2$
  \RightLabel{\Intro\lif}
  \UnaryInf$\fCenter \lambd[\typeof{x}{!A \land
  !B}][\pair{\proj{2}{x}}{\proj{1}{x}}] : (!A \land !B) \lif (!C_1 \land !C_2)$
  \DisplayProof
  \]
ترم برهان~$\proj{2}{x}$ نشان می‌دهد دنبالهٔ بالایی سمت چپ با
$\Elim{\land}[2]$ به دست آمده و بنابراین مقدمه به شکل~$!D \land !C_1$
است. در سمت راست می‌توان تعیین کرد استنتاج منتهی به~$!C_2$ باید
$\Elim{\land}[1]$ باشد و مقدمه به شکل~$!C_2 \land !E$ باشد.
  \[
  \Axiom$x : !A\land !B \fCenter x: !D\land !C_1$
  \RightLabel{\Elim\land[2]}
  \UnaryInf$x : !A\land !B \fCenter \proj{2}{x} : !C_1$
  \Axiom$x : !A\land !B \fCenter x: !C_2\land !E$
  \RightLabel{\Elim\land[1]}
  \UnaryInf$x : !A\land !B \fCenter \proj{1}{x} : !C_2$
  \RightLabel{\Intro\land}
  \BinaryInf$x : !A\land !B \fCenter \pair{\proj{2}{x}}{\proj{1}{x}} :
    !C_1 \land !C_2$
  \RightLabel{\Intro\lif}
  \UnaryInf$\fCenter \lambd[\typeof{x}{!A \land
  !B}][\pair{\proj{2}{x}}{\proj{1}{x}}] : (!A \land !B) \lif (!C_1 \land !C_2)$
  \DisplayProof
  \]
چون ترم‌های برهان دنباله‌های بالایی اکنون متغیرند، می‌دانیم آن
دنباله‌ها باید اصل باشند. این امر~$!C_1$، $!C_2$، $!D$ و~$!E$ را
تعیین می‌کند: باید~$!D \ident !A$ و~$!C_1 \ident !B$ باشد، وگرنه
دنبالهٔ چپ اصل نیست؛ و باید~$!C_2 \ident !A$ و~$!E \ident !B$ باشد،
وگرنه دنبالهٔ راست اصل نیست. اکنون برهان را کاملاً بازیابی کرده‌ایم:
  \[
  \Axiom$x : !A\land !B \fCenter x: !A\land !B$
  \RightLabel{\Elim\land[2]}
  \UnaryInf$x : !A\land !B \fCenter \proj{2}{x} : !B$
  \Axiom$x : !A\land !B \fCenter x: !A\land !B$
  \RightLabel{\Elim\land[1]}
  \UnaryInf$x : !A\land !B \fCenter \proj{1}{x} : !A$
  \RightLabel{\Intro\land}
  \BinaryInf$x : !A\land !B \fCenter \pair{\proj{2}{x}}{\proj{1}{x}} :
    !B \land !A$
  \RightLabel{\Intro\lif}
  \UnaryInf$\fCenter \lambd[\typeof{x}{!A \land
  !B}][\pair{\proj{2}{x}}{\proj{1}{x}}] : (!A \land !B) \lif (!B \land !A)$
  \DisplayProof
  \]


\end{document}
